% Entrain: draft skeleton. Every number is generated from bench/results.json
% (paper/tables/*.tex, see paper/make_tables.py); every formula is transcribed
% from entrain_infer/theorem.py and entrain_infer/detector.py.
%
% Status marks: \proved (proof in the text or in theorem.py), \measured (read
% from bench/results.json), \assumed (an assumption, not a result).
\documentclass[10pt,twocolumn,a4paper]{article}

\usepackage[margin=1.9cm]{geometry}
\usepackage{amsmath,amssymb,amsthm}
\usepackage{booktabs}
\usepackage{graphicx}
% microtype needs scalable EC fonts, which this TeX installation lacks; the
% draft compiles with plain Computer Modern Type1 (amsfonts) instead.
% \usepackage{microtype}
\usepackage{xcolor}
\usepackage[colorlinks=true,allcolors=blue!55!black]{hyperref}

% ── status marks ────────────────────────────────────────────────────────────
\newcommand{\proved}{\textcolor{teal!70!black}{\textsc{[proved]}}}
\newcommand{\measured}{\textcolor{blue!70!black}{\textsc{[measured]}}}
\newcommand{\assumed}{\textcolor{orange!80!black}{\textsc{[assumed]}}}
\newcommand{\gap}[1]{\textcolor{red!70!black}{\textbf{[to write: #1]}}}

\newtheorem{theorem}{Theorem}[section]
\newtheorem{lemma}[theorem]{Lemma}
\newtheorem{corollary}[theorem]{Corollary}
\newtheorem{proposition}[theorem]{Proposition}
\theoremstyle{definition}
\newtheorem{definition}[theorem]{Definition}
\newtheorem{assumption}[theorem]{Assumption}
\theoremstyle{remark}
\newtheorem{remark}[theorem]{Remark}

\newcommand{\PH}{P_{H}}
\newcommand{\one}{\mathbf{1}}
\newcommand{\R}{\mathbb{R}}
\newcommand{\bmin}{b_{\min}}
\newcommand{\taur}{\tau}

\title{\bfseries Entrain: a gauge symmetry, a chord, and a one-sided bound\\
\large What a park of oscillators can and cannot say about a lying clock}

\author{ENRG project\thanks{Correspondence: \texttt{enrg.project@gmail.com}.
Author list and affiliations to be finalized \gap{authors, affiliation, ORCID}.}}
\date{Draft --- \today}

\begin{document}
\maketitle

\begin{abstract}
A park of $N$ oscillators keeps time without an external clock: every node reads
its own frequency, and the only thing the park can compare is the nodes against
each other. We ask what such a park can prove about a node that lies, and what it
provably cannot. The detector is a generalized likelihood ratio test on
studentized contrasts, with its threshold taken from a declared false-alarm risk
$\alpha$ rather than from a tuned constant; on a measured background its power at
$\alpha=10^{-3}$ matches the closed form to the third decimal and its ROC area is
$0.991$--$0.993$ \measured. What follows is a statement about the observation, not
about the test: the park sees the quotient $\R^N/\mathrm{span}\{\one\}$, so a lie
enters only through its contrast $\PH\one_S$, whose length is the chord
$\sqrt{k(N-k)/N}$ \proved. Three consequences. A coalition of all $N$ nodes is
invisible at any bias and any power; a coalition of $k>N/2$ is detected while the
honest minority is named; the orientation of a split -- who lies, the $k$ or the
$N-k$ -- is not estimable, because the two descriptions are the same vector.
Dilution obeys a ceiling $\bmin(k)\le\bmin(1)\,(N-1)/\max(k,N-k)$ with equality
only at the two ends \proved, and the only converse is an anchor, which buys
immunity to large coalitions at the price of a $2.4\times$ ($N=10$) to $4.3\times$
($N=30$) worse floor against a lone liar \measured. Every claim is marked as
proved, measured, or assumed, and we name the places where the model, not the
mathematics, is the weak link.
\end{abstract}

\section{Introduction}
\label{sec:intro}

Distributed systems that keep time by consensus have a structural problem: the
thing they most need to verify -- the clock -- is the one thing no member can
check against anything external. A solar inverter, a substation relay, a sensor
in a park: each reports a frequency, and the report is meaningful only relative
to the other reports. If one node drifts, the park sees that \emph{something}
moved. If one node lies, the park sees exactly the same thing. The question of
this paper is not how to detect drift, but what a park can \emph{prove} about a
lie when the only observable is the park itself.

We take the classical route and make it honest about its assumptions. The
observable is a matrix of per-node frequencies over a window; the test is a GLRT
on studentized contrasts, in the lineage of \cite{lehmann2005}; the threshold is
a quantile of the null at a declared risk $\alpha$, computed once per $N$ and
printed next to the Bonferroni bound, so that an operator can read the
false-alarm price in the same units in which the detection claim is made
\proved. This engineering half is measured rather than asserted: on real
background traces the achieved false-positive rate tracks $\alpha$ and the
achieved power tracks the closed form to the third decimal
(\S\ref{sec:glrt}, Table~\ref{tab:roc}).

The mathematical half is different in kind, and it is the reason for this paper.
Once the detector is written as a function of the contrasts, a symmetry appears
that no amount of data can break: adding the same number to every node moves
nothing. This is not an artefact of studentization; it says that the observation
lives in a quotient space of dimension $N-1$, and every question the park can
answer is a question about that quotient. In it:

\begin{enumerate}
\item A coalition of all $N$ nodes agreeing on one bias is \emph{invisible} at
      any bias, any window, any power. There is no threshold to lower and no
      sample to collect: alternative and null coincide
      (Theorem~\ref{thm:kernel}).
\item The visibility of a coalition is a chord,
      $c(S)=\|\PH\one_S\|=\sqrt{k(N-k)/N}$: zero at the ends, maximal at $N/2$
      (Lemma~\ref{lem:chord}). Detectability is a property of the \emph{split},
      not of the lie.
\item For $k>N/2$ the largest statistics belong to the honest minority: the
      detector fires correctly and names the wrong nodes
      (\S\ref{sec:inversion}). This is a statement about the statistic,
      confirmed on the measured background, not an implementation bug.
\item The orientation of a split -- ``$k$ lied by $+b$'' versus ``$N-k$ lied by
      $-b$ while the park clock is off by $+b$'' -- is not estimable, because the
      two worlds are one vector $x$. A majority vote is a choice of gauge, not an
      inference (Lemma~\ref{lem:orientation}).
\item Dilution obeys a ceiling, and the ceiling is reached only at the two ends;
      between them the bound is strict, and we prove the direction
      (Lemma~\ref{lem:bound}).
\item The only converse is to pin the gauge by hand: declare one node an anchor.
      The anchored floor falls with $k$ while the relative one grows, so the
      anchor is a trade, worth its assumption only against coalitions above
      $k^\ast\approx\sqrt{N}$ (\S\ref{sec:anchor}).
\end{enumerate}

None of the six is a claim about a better algorithm. They are claims about the
information in the observation, and they bind any algorithm that respects the
gauge, including ones not yet written. That is what makes them worth stating, and
why they come with a list of what the model does not cover
(\S\ref{sec:limits}).

\medskip
\noindent\textbf{What is proved, what is measured, what is assumed.} Claims are
marked throughout. \proved{} means a proof in the text or a closed-form identity
in \texttt{theorem.py}, with no Monte Carlo and no seed. \measured{} means a
number from the frozen run \texttt{bench/results.json}, reproducible to the byte
(Appendix~\ref{app:repro}). \assumed{} means an assumption of the model: Gaussian
stationary noise, a single bias $b$ for the whole coalition, a clock that drifts
rather than jumps. The distinction matters -- the symmetry results survive any
noise model, the quantitative ones do not.

\section{Model}
\label{sec:model}

\begin{definition}[Park]
\label{def:park}
A park of $N\ge2$ oscillators is a vector $x\in\R^N$, $x_i$ the frequency of
node $i$ read over a window of $T$ samples. Under $H_0$ the nodes share a common
drift $\mu$ and independent noise: $x=\mu\one+\sigma z$, $z_i\sim\mathcal N(0,1)$.
Under $H_1$ a set $S\subseteq\{1,\dots,N\}$ of $k=|S|$ nodes carries one bias,
\begin{equation}
x=\mu\one+b\,\one_S+\sigma z .
\label{eq:h1}
\end{equation}
\end{definition}

Three things about \eqref{eq:h1} are assumptions and nothing more \assumed. The
noise is Gaussian and stationary over the window: the measured background is not
strictly so, and \S\ref{sec:field} reports the excess kurtosis left after
detrending, up to $1.689$ in one cell \measured. The bias is a single number for
the whole coalition: a coalition split into two opposite biases is a different,
weaker model, and \S\ref{sec:anchor} shows why. The lie is a drift, not a jump: a
timestamp step is a common-mode shift and is invisible to a frequency reading for
the same reason as Theorem~\ref{thm:kernel}.

\begin{remark}[What the park observes]
The vector $x$ is never available in absolute terms: only differences between
nodes carry information about $\mu$ or $b$. Formally the observation is the image
of $x$ under the orthogonal projection
\begin{equation}
\PH v = v-\bar v\,\one,\qquad \bar v=\frac1N\sum_j v_j,
\label{eq:ph}
\end{equation}
i.e.\ a point of the hyperplane $H=\one^{\perp}\cong\R^{N-1}$. The statistic of
\S\ref{sec:glrt} is a function of $\PH x$ alone, and so is every statistic that
respects the gauge; \S\ref{sec:theorem} makes this precise.
\end{remark}

\section{The GLRT detector}
\label{sec:glrt}

The test is the classical one. Studentize the contrasts of \eqref{eq:ph},
\begin{equation}
t_i=\frac{N x_i-\sum_j x_j}{\sqrt{N(N-1)}}=\frac{(\PH x)_i}{\|\PH e_i\|},
\label{eq:stat}
\end{equation}
and reject $H_0$ for a node when $|t_i|$ exceeds $\taur$. Under $H_0$ the vector
$t$ is exchangeable with unit variance and pairwise correlation $-1/(N-1)$; under
$H_1$ with a lone liar the offending node is noncentral,
\begin{equation}
\mathbb E\,t_i=\frac{b}{\sqrt{N-1}}\quad (i\in S),\qquad
\mathbb E\,t_i=0\quad\text{otherwise.}
\label{eq:tmean1}
\end{equation}

\paragraph{Threshold as a declared risk.} $\taur$ is the $(1-\alpha)$ quantile of
$\max_i|t_i|$ under $H_0$, computed by Monte Carlo once per $N$
(200\,000 null draws, seed 20261030) and printed next to the Bonferroni value
$\Phi^{-1}(1-\alpha/2N)$; the run checks the two against each other. At
$\alpha=10^{-3}$ the measured thresholds are $\taur=3.875$, $4.154$, $4.424$ for
$N=10,30,100$ \measured. The point is not the number but its units: $\alpha$ is
the share of honest parks that will be flagged, so the operator pays a stated
price instead of inheriting a constant. Table~\ref{tab:roc} shows the achieved
false-positive rate tracking $\alpha$ and the achieved power matching the closed
form.

\paragraph{Minimum detectable lie.} Because \eqref{eq:stat} is scale invariant,
$\bmin$ is linear in $\sigma$:
\begin{equation}
\bmin(N,\sigma)=\sigma\,c(N),\qquad c(N)=\bmin(N,\sigma=1),
\label{eq:bmin}
\end{equation}
with $c(10)=5.420$, $c(30)=5.513$, $c(100)=5.740$ \measured{} at $\alpha=10^{-3}$
and target power $0.9$. One coefficient per $N$ therefore serves every season,
because the test is a function of $b/\sigma$. The closed-form companion,
$\bmin^{\mathrm{an}}=({\taur}+z_{1-\beta})\,\sigma\sqrt{1+1/(N-1)}$, agrees with
the bisected value to $0.3\%$; the gap is the difference between the Monte-Carlo
quantile and the bisection on Monte-Carlo power \measured.

\paragraph{ROC.} Table~\ref{tab:roc} gives four operating points per $N$ at
$T=24$~h on the July background and the area under the measured ROC curve:
$0.993$, $0.991$, $0.991$ for $N=10,30,100$ \measured. Two regularities are worth
naming. The curve is flat in $T$ -- the statistic depends on $b/\sigma$ and on
$N$, while $T$ enters only through the estimate of $\sigma$ -- and it degrades
slowly in $N$ at fixed $b/\sigma$, because the maximum of $N$ correlated unit
normals grows like $\sqrt{2\log N}$ while the noncentrality of one liar grows
like $\sqrt{N-1}$. A larger park does not help against a single liar;
\S\ref{sec:anchor} shows the one setting in which it does.

\begin{table*}[t]
\centering
\caption{ROC of the declared risk (\texttt{roc}, July, $T=24$~h, frequency offset). $\tau$ is the Monte-Carlo quantile of the null (200\,000 draws, seed 20261030, standard error printed in the report); the achieved false-positive rate tracks $\alpha$ and the achieved power matches the closed-form value to the third decimal. The threshold is not a constant but a function of the declared risk and of $N$ -- which is why a magic $0.5$ is replaced by $\alpha$ in the operator's units.}
\label{tab:roc}
\footnotesize
\setlength{\tabcolsep}{4pt}
\begin{tabular}{@{}l r r r r r@{}}
\toprule
 & $\alpha$ & $\tau$ & achieved FPR & achieved TPR & analytic TPR \\
\midrule
\multicolumn{6}{@{}l}{$N=100$, AUC $=0.991$} \\
 & $0.001$ & 4.4238 & 0.001050 & 0.9022 & 0.9010 \\
 & $0.003$ & 4.1809 & 0.003200 & 0.9365 & 0.9370 \\
 & $0.01$ & 3.8967 & 0.009875 & 0.9655 & 0.9652 \\
 & $0.1$ & 3.2787 & 0.100150 & 0.9935 & 0.9925 \\
\multicolumn{6}{@{}l}{$N=10$, AUC $=0.993$} \\
 & $0.001$ & 3.8753 & 0.000800 & 0.8975 & 0.8973 \\
 & $0.003$ & 3.6095 & 0.002750 & 0.9378 & 0.9373 \\
 & $0.01$ & 3.2912 & 0.010100 & 0.9695 & 0.9679 \\
 & $0.1$ & 2.5521 & 0.101750 & 0.9959 & 0.9952 \\
\multicolumn{6}{@{}l}{$N=30$, AUC $=0.991$} \\
 & $0.001$ & 4.1544 & 0.001025 & 0.8958 & 0.8972 \\
 & $0.003$ & 3.8861 & 0.003225 & 0.9380 & 0.9375 \\
 & $0.01$ & 3.5875 & 0.009725 & 0.9666 & 0.9666 \\
 & $0.1$ & 2.9185 & 0.100675 & 0.9943 & 0.9938 \\
\bottomrule
\end{tabular}
\end{table*}


\gap{implementation paragraph: the estimator of the common drift and of $\sigma$
(window, night fraction $0.15$, one clock per park) and the reason
\eqref{eq:stat} does not depend on the estimate of $\mu$.}

\section{The gauge symmetry}
\label{sec:theorem}

\begin{lemma}[Exchange]
\label{lem:gauge}
For every $c\in\R$ and every $x\in\R^N$, $t(x+c\one)=t(x)$; equivalently, $t$ is
a function of $\PH x$ alone, and $\PH$ is the orthogonal projector onto
$\one^{\perp}$.
\end{lemma}

\begin{proof}
$t$ is built from $x_i-\bar x$, and $\bar x$ absorbs $c$; in matrix form
$t=\PH x/\sqrt{\PH}$-normalized, so $t(x+c\one)=t(x)$. Since $\PH$ is symmetric
and idempotent, $\PH$ is the orthogonal projector onto $\mathrm{span}\{\one\}^\perp$;
its rank is $N-1$ and its kernel is $\mathrm{span}\{\one\}$.
\end{proof}

The lemma is trivial, and that is the point: the triviality is a symmetry, not a
technicality. The observation is a point of a quotient space, so a lie that
differs from the null only by a vector in the kernel is not a lie the park can
see -- not a hard detection problem, an empty one.

\begin{theorem}[The kernel is silent]
\label{thm:kernel}
If $S=\{1,\dots,N\}$ (every node lies by the same $b$), then
$\PH\one_S=\PH\one=0$, so $x$ of \eqref{eq:h1} and $x$ of $H_0$ with
$\mu\to\mu+b$ are the same observation, $t\equiv0$, and for every measurable
rejection region with level $\alpha$ the power is exactly $\alpha$.
\end{theorem}

\begin{proof}
$\one_S=\one\in\ker\PH$, so $\PH\one_S=0$ and $\PH x=\PH(\mu\one+b\one+\sigma z)=\sigma\PH z$.
The two hypotheses induce the same law of $\PH x$; by the Neyman--Pearson lemma
\cite{neyman1933} no test has power above its size.
\end{proof}

Table~\ref{tab:kernel} is the measured shadow of the theorem: on the real
background with every node lying at $b/\sigma=1,5,20$, the recovered bias is
$\hat b_{\text{eff}}/\sigma\sim10^{-17}$ and detection stays at $0.000$
\measured. Increasing $b$ cannot help, because the theorem says the direction of
increase is in the kernel.

\begin{table}[t]
\centering
\caption{The kernel is a silence, not a weak signal (\texttt{kernel\_real}). Every node lies by the same $b$: the statistic is zero up to floating point ($\hat b_{\text{eff}}/\sigma\sim10^{-17}$) at $b/\sigma=1,5,20$, so no increase in $b$ can ever buy detection. Measured detection is $0.000$ in all three cells, against the $\alpha=0.001$ a coin would give.}
\label{tab:kernel}
\small
\begin{tabular}{@{}r r r r@{}}
\toprule
$b/\sigma$ & $\hat b_{\text{eff}}/\sigma$ & detection & localisation \\
\midrule
1.0 & -0.000 & 0.000 & 1.000 \\
5.0 & -0.000 & 0.000 & 1.000 \\
20.0 & -0.000 & 0.000 & 1.000 \\
\bottomrule
\end{tabular}
\end{table}



\begin{remark}[What the theorem does \emph{not} say]
It does not say that a lying node is undetectable -- only that a
\emph{unanimous} lie is. It does not depend on the noise being Gaussian: any
law invariant to common-mode shifts gives the same conclusion. And it does not
depend on $t$ being the GLRT: any statistic measurable with respect to $\PH x$
obeys it. This is why we call it a statement about the observation.
\end{remark}

\section{The chord: how visible a split is}
\label{sec:chord}

Theorem~\ref{thm:kernel} is the degenerate end of a family. What the park sees of
a coalition is the projection of its indicator, and the length of that projection
is the only thing about the coalition that enters the data.

\begin{lemma}[Chord]
\label{lem:chord}
For $S$ with $|S|=k$,
\begin{equation}
\begin{aligned}
c(S)^2&=\|\PH\one_S\|^2=\frac{k(N-k)}{N},\\
c(S)&=\sqrt{\frac{k(N-k)}{N}} .
\end{aligned}
\label{eq:chord}
\end{equation}
$c$ is symmetric under $k\mapsto N-k$, vanishes at $k=0$ and $k=N$, and is
maximal at $k=N/2$.
\end{lemma}

\begin{proof}
$\|\one_S\|^2=k$ and $\langle\one_S,\one\rangle=k$, so
$\|\PH\one_S\|^2=\|\one_S\|^2-\frac1N\langle\one_S,\one\rangle^2=k-k^2/N$.
The rest is calculus on $k(N-k)$.
\end{proof}

The consequence is a reparametrisation of the detection problem: a coalition of
$k$ liars lying by $b$ is the same observation as one liar lying by
$b\,c(S)/c(\{1\})=b\sqrt{k(N-k)/(N-1)}$, up to the second-order question of how
the test is ordered (\S\ref{sec:bound}). Detectability grows with the
\emph{product} $k(N-k)$, not with $k$: a coalition of nine in a park of ten is as
visible as a single liar, and a coalition of five is the most visible thing that
can happen in a park of ten (Table~\ref{tab:chord}).

\begin{proposition}[Per-node means]
\label{prop:means}
Under \eqref{eq:h1}, for $i\in S$ and $j\notin S$,
\begin{equation}
\mathbb E\,t_i=\frac{b\,(N-k)}{\sqrt{N(N-1)}},\qquad
\mathbb E\,t_j=-\frac{b\,k}{\sqrt{N(N-1)}} .
\label{eq:means}
\end{equation}
\end{proposition}

\begin{proof}
$(\PH\one_S)_i=1-k/N=(N-k)/N$ for $i\in S$ and $-k/N$ for $j\notin S$; divide by
$\|\PH e_i\|=\sqrt{(N-1)/N}$.
\end{proof}

Equation~\eqref{eq:means} contains M1's dilution law and its inversion in one
line. For $k\le N/2$ the colluders carry the larger mean and the detector's
answer is correct up to the sign; the measured power of a lone liar is recovered
at $k=N-1$ (Table~\ref{tab:measured}, $N=10$, $k=9$: power $0.897$ against
$0.899$ for $k=1$) \measured. For $k>N/2$ the honest minority carries the larger
mean and the detector names them -- the subject of \S\ref{sec:inversion}.

\begin{table*}[t]
\centering
\caption{The chord and the per-node statistic, closed form (\texttt{theorem.predicted}); no Monte Carlo in this table. $c(S)=\|P_H\mathbf{1}_S\|=\sqrt{k(N-k)/N}$, so $c/c_1$ is the detectability of a known coalition relative to one liar; $\mathbb{E}t_i/b$ are the per-node means inside and outside the coalition; gain is $\max(k,N-k)/(N-1)$, the noncentrality of the largest $|t|$ relative to a single liar. ``flagged'' is the side the GLRT names. At $k=N$ every $t_i$ is identically zero: nothing is named, because there is no second side to compare against.}
\label{tab:chord}
\small
\begin{tabular}{@{}r r r r r r r l@{}}
\toprule
$N$ & $k$ & $c/c_1$ & in & out & gain & $b_{\min}$ ratio & flagged \\
\midrule
10 & 1 & 1.000 & 0.949 & -0.105 & 1.000 & 1.000 & colluder \\
10 & 2 & 1.333 & 0.843 & -0.211 & 0.889 & 1.125 & colluder \\
10 & 3 & 1.528 & 0.738 & -0.316 & 0.778 & 1.286 & colluder \\
10 & 5 & 1.667 & 0.527 & -0.527 & 0.556 & 1.800 & tie \\
10 & 6 & 1.633 & 0.422 & -0.632 & 0.667 & 1.500 & honest \\
10 & 8 & 1.333 & 0.211 & -0.843 & 0.889 & 1.125 & honest \\
10 & 9 & 1.000 & 0.105 & -0.949 & 1.000 & 1.000 & honest \\
10 & 10 & 0.000 & 0.000 & -1.054 & 0.000 & $\infty$ & none \\
30 & 1 & 1.000 & 0.983 & -0.034 & 1.000 & 1.000 & colluder \\
30 & 2 & 1.390 & 0.949 & -0.068 & 0.966 & 1.036 & colluder \\
30 & 5 & 2.076 & 0.848 & -0.170 & 0.862 & 1.160 & colluder \\
30 & 10 & 2.626 & 0.678 & -0.339 & 0.690 & 1.450 & colluder \\
30 & 15 & 2.785 & 0.509 & -0.509 & 0.517 & 1.933 & tie \\
30 & 20 & 2.626 & 0.339 & -0.678 & 0.690 & 1.450 & honest \\
30 & 25 & 2.076 & 0.170 & -0.848 & 0.862 & 1.160 & honest \\
30 & 29 & 1.000 & 0.034 & -0.983 & 1.000 & 1.000 & honest \\
30 & 30 & 0.000 & 0.000 & -1.017 & 0.000 & $\infty$ & none \\
\bottomrule
\end{tabular}
\end{table*}


\begin{remark}[The chord test]
\label{rem:chordtest}
If the coalition is known in advance -- an operator's own devices, say -- the
efficient test is not the node maximum but the two-cluster contrast
\begin{equation}
z=\frac{\bar x_S-\bar x_{S^c}}{\sigma\sqrt{1/k+1/(N-k)}}
\ \sim\ \mathcal N\!\left(\frac{b\,c(S)}{\sigma},\,1\right),
\label{eq:chordtest}
\end{equation}
whose noncentrality is exactly the chord. Its minimum detectable lie is
$\bmin^{\mathrm{ch}}=({\taur}+z_{1-\beta})\sigma/c(S)$, finite for every $k<N$
and infinite at $k=N$ -- Theorem~\ref{thm:kernel} as a division by zero. The gap
between \eqref{eq:chordtest} and $\max_i|t_i|$ is a \emph{design} gap, not an
information gap: both are functions of $\PH x$, but one is ordered by the
statistic that carries the chord, the other by the largest coordinate. The run
measures both side by side: at $k=N/2$ the chord power at $b_1$ is $1.000$ where
the node maximum gives $0.793$ ($N=10$) \measured.
\end{remark}

\section{Inversion and orientation}
\label{sec:inversion}

\begin{corollary}[Inversion]
\label{cor:inversion}
By \eqref{eq:means}, $\mathbb E\,t_i<\mathbb E\,t_j$ for $i\in S$, $j\notin S$
whenever $k>N/2$: the largest statistics of the park belong to the honest
minority, and a detector that flags the largest $|t|$ names the honest.
\end{corollary}

This is where the engineering reading and the mathematical reading part company.
An operator sees a park in which the detector fired and the flagged nodes are the
ones that moved. The mathematics says the flagged nodes are $N-k$, the smaller
side, and that the side with $k$ is the one that moved. Both statements are true
of the same numbers; only the second is about the data.

The measured background does not soften this. In Table~\ref{tab:field}, past
$k>N/2$ detection stays at $1.000$ while localisation -- the share of repetitions
whose largest $|t|$ sits on a colluder -- falls to $0.000$: at $N=10$, $k=9$, the
detector fires in $100\%$ of repetitions and names an honest node in $100\%$ of
them \measured. A confident detector, right about the existence of a split and
wrong about its direction.

\begin{lemma}[Orientation is in the kernel]
\label{lem:orientation}
For every $S$, $b$ and $\mu$,
\begin{equation}
\mu\one+b\one_S=(\mu+b)\one-b\one_{S^c}
\label{eq:orientation}
\end{equation}
identically, as vectors in $\R^N$. The two worlds -- ``$k$ nodes lie by $+b$''
and ``$N-k$ nodes lie by $-b$ while the park clock is off by $+b$'' -- induce the
same law of the observation, hence the same $t$, the same flagged node, and the
same $\bmin$.
\end{lemma}

\begin{proof}
$\one_S+\one_{S^c}=\one$; expand the right-hand side.
\end{proof}

So the orientation of a split is not a hard estimation problem: it is not in the
data. Every statement of the form ``the majority is honest'' is a choice of gauge,
and every vote is a convention for fixing it. This is the same freedom that makes
a constant timestamp offset invisible to a frequency reading; there the gauge is
fixed by an external time reference, here by a decision about whom to believe. The
run checks \eqref{eq:orientation} on 15 cells: apart from floating point, the
reading gap and the statistic gap are zero and the flagged node is the same in
both worlds \proved.

\begin{remark}[What a vote can and cannot buy]
A vote among nodes is a gauge choice, and a gauge choice can be wrong. What a vote
\emph{can} do is define $S$ well enough for the chord test of
Remark~\ref{rem:chordtest} to be run with roles assigned: given a candidate
coalition, the statistic says whether the park is split along it, not which side
is lying. Anything stronger needs information from outside the park -- an anchor
(\S\ref{sec:anchor}), a trusted time reference, or a physical model of who can
drift in which direction.
\end{remark}

\section{The anchor as a trade}
\label{sec:anchor}

The only way out is to stop asking a question the observation cannot answer.
Declare one node's clock true; then
\begin{equation}
y_i=x_i-x_a,\qquad i\neq a,
\label{eq:anchor}
\end{equation}
has a \emph{known} mean under $H_0$, and the question stops being ``which node
deviates from the others'' (a relative question, $N-1$ contrasts) and becomes
``does the rest deviate from the anchor'' (an absolute one).

\begin{proposition}[Anchored test: the i.i.d.\ convention]
\label{prop:anchor}
Let $y\in\R^{N-1}$ be the standardised anchored readout, $y_i\sim\mathcal N(0,\sigma^2)$
independent, and $s=\sqrt{N-1}\,\bar y/\sigma$. Under $H_0$, $s\sim\mathcal N(0,1)$
exactly: no null table, no $N$-dependent quantile, because the mean under $H_0$ is
known rather than profiled. With $k$ liars of the same sign past the anchor,
\begin{equation}
\delta=\frac{b\,k}{\sqrt{N-1}},\qquad
\bmin^{\mathrm{anc}}=\frac{({\taur}+z_{1-\beta})\,\sigma\sqrt{N-1}}{k},
\label{eq:anchored}
\end{equation}
finite for every $k$ -- including $k=N-1$, where the relative test is as
insensitive as against a lone liar \emph{and} names the honest node, while the
anchor is at its best, $\bmin^{\mathrm{anc}}=({\taur}+z_{1-\beta})\sigma/\sqrt{N-1}$.
\end{proposition}

\begin{proof}
Under $H_0$ each $y_i$ is centred, so $\bar y$ has mean $0$ and variance
$\sigma^2/(N-1)$; $\sqrt{N-1}\,\bar y/\sigma$ is standard normal. Under $H_1$ with
$k$ of the $N-1$ coordinates shifted by $b$, $\mathbb E\bar y=bk/(N-1)$, hence
$\delta=bk/\sqrt{N-1}$. Inverting $1-\Phi(\taur-\delta)=1-\beta$ gives
\eqref{eq:anchored}.
\end{proof}

\begin{remark}[Scaling convention, and a discrepancy we did not hide]
\label{rem:scaling}
\eqref{eq:anchored} presumes $y$ is a readout whose own noise is $\sigma$ --
the convention of the bench, which draws $y$ directly and adds the lie.
If instead $y_i=x_i-x_a$ is built from the park of \eqref{eq:h1}, the coordinates
are correlated: $\operatorname{Var}(y_i)=2\sigma^2$ and
$\operatorname{Cov}(y_i,y_j)=\sigma^2$, so
$\operatorname{Var}(\bar y)=\sigma^2(1/(N-1)+1)=\sigma^2N/(N-1)$ and the statistic
$\sqrt{N-1}\,\bar y/\sigma$ has variance $N$, not $1$ \proved. The correctly normalised readout is
$\bar y\sqrt{(N-1)/N}/\sigma$, and with $\sigma_p\equiv\sigma_{\mathrm{park}}$ the
anchored noncentrality becomes
\begin{equation}
\begin{aligned}
\delta_{\mathrm{park}}&=\frac{b\,k}{\sigma_p\sqrt{N(N-1)}},\\
\bmin^{\mathrm{anc,park}}&=\frac{({\taur}+z_{1-\beta})\,\sigma_p\sqrt{N(N-1)}}{k},
\end{aligned}
\label{eq:anchored-park}
\end{equation}
a factor $\sqrt{N}$ worse than \eqref{eq:anchored} at the same $b$, while still
linear in $1/k$ and still finite at $k=N-1$ \proved. The crossover $k^\ast$, the
sign of the trade, and the ordering of the two tests are unchanged; the
\emph{price} of the anchor against a lone liar is larger than
Table~\ref{tab:anchor} reports. \gap{decide the convention, then make
\texttt{theorem.anchor\_statistic} agree with it (it currently returns a statistic
of variance $N\sigma^2$ on park data) and regenerate Table~\ref{tab:anchor} in
both conventions.}
\end{remark}

The anchored floor falls with $k$ where the relative one grows, and the anchored
test needs no multiplicity correction (one statistic, not $N$). The two curves
cross once, and where they cross is the whole governance question.

\begin{table}[t]
\centering
\caption{The anchor as a trade (\texttt{theorem.anchor}). An anchored test compares $x_i-x_a$, whose mean is known, so $b_{\min}=(\tau+z)\sigma\sqrt{N-1}/k$ falls with $k$ while the relative floor grows up to $N/2$. ``gain'' is the relative floor over the anchored one: below $k^\ast$ it is $<1$, i.e.\ against a few liars the anchor is worse than the plain GLRT it replaces. $k^\ast=3$ at $n=10$ ($\sqrt{n}=3.162$), $k^\ast=4$ at $n=30$ ($\sqrt{n}=5.477$). At $k=N$ the anchor itself lies, $y\equiv0$, and the measured power falls back to $\alpha$ (0.0013 at $N=30$).}
\label{tab:anchor}
\small
\setlength{\tabcolsep}{4pt}
\begin{tabular}{@{}r r r r r r@{}}
\toprule
$N$ & $k$ & anchor $b_{\min}/\sigma$ & relative $b_{\min}/\sigma$ & gain & power \\
\midrule
10 & 1 & 13.115 & 5.420 & 0.413 & 0.8981 \\
10 & 2 & 6.558 & 6.097 & 0.930 & 0.8981 \\
10 & 3 & 4.372 & 6.968 & 1.594 & 0.8981 \\
10 & 5 & 2.623 & 9.756 & 3.719 & 0.8981 \\
10 & 6 & 2.186 & 8.130 & 3.719 & 0.8981 \\
10 & 8 & 1.639 & 6.097 & 3.719 & 0.8981 \\
10 & 9 & 1.457 & 5.420 & 3.719 & 0.8981 \\
10 & 10 & $\infty$ & $\infty$ & -- & 0.0017 \\
30 & 1 & 23.543 & 5.513 & 0.234 & 0.9013 \\
30 & 2 & 11.771 & 5.710 & 0.485 & 0.9013 \\
30 & 5 & 4.709 & 6.395 & 1.358 & 0.9013 \\
30 & 10 & 2.354 & 7.993 & 3.395 & 0.9013 \\
30 & 15 & 1.570 & 10.658 & 6.791 & 0.9013 \\
30 & 20 & 1.177 & 7.993 & 6.791 & 0.9013 \\
30 & 25 & 0.942 & 6.395 & 6.791 & 0.9013 \\
30 & 29 & 0.812 & 5.513 & 6.791 & 0.9013 \\
30 & 30 & $\infty$ & $\infty$ & -- & 0.0013 \\
\bottomrule
\end{tabular}
\end{table}



\section{The one-sided dilution bound}
\label{sec:bound}

The chord is the information; the node-maximum statistic is what M1 uses, and the
two are ordered differently. What the run measured for the node maximum -- and
what M1 reported as a dilution law -- needs a proof of its direction, and that is
this section.

\begin{lemma}[The bound is a ceiling]
\label{lem:bound}
Let $m=\max(k,N-k)$, $s=\min(k,N-k)$ and $r=(N-1)/m$. Then
\begin{equation}
\frac{\bmin(k)}{\bmin(1)}\ \le\ r=\frac{N-1}{\max(k,N-k)},
\label{eq:bound}
\end{equation}
with equality if and only if $k\in\{1,N-1\}$; for $2\le k\le N-2$ the inequality
is strict.
\end{lemma}

\begin{proof}
Take $k\le N/2$ without loss: the map $k\mapsto N-k$ is symmetric, because
$c(S)=c(S^c)$ by Lemma~\ref{lem:chord}. Raise the coalition's lie to
$b'=b\,r$. Its $k$ largest statistics sit at the noncentrality a lone liar has at
$b$:
\begin{equation}
\mathbb E\,t_i=b'\,\frac{N-k}{\sqrt{N(N-1)}}=b\sqrt{\frac{N-1}{N}},
\label{eq:lemma-mean}
\end{equation}
while the $N-k$ nodes on the other side sit at $-b'k/\sqrt{N(N-1)}$. The threshold
$\taur$ is the same in both worlds, so for the target power of one liar
\begin{equation}
1-(1-p)^{m}\le\mathrm{power}(\text{coalition at }b')
\iff \bmin(k)\le b\,r ,
\label{eq:lemma-power}
\end{equation}
where $p$ is the chance that one node at that noncentrality crosses $\taur$ and
$m$ the number of nodes sharing it. Equality needs $m=1$, i.e.\ $k=N-1$; at
$k=1$ we have $r=1$ and the configuration is the same one on both sides. For
$2\le k\le N-2$ we have $m\ge2$ and
$(1-p)^{m-1}<(1-\alpha)^{N-1}$: $m$ nodes sharing the maximal noncentrality fire
more often than the one node that alone sets $\bmin(1)$, so the inequality is
strict.
\end{proof}

The lemma is stated one-sidedly on purpose, and the run reports the slack as a
column instead of hiding it behind a tolerance. Table~\ref{tab:measured} shows
both: the predicted ceiling, the measured ratio, and a ``bound'' column that is
\texttt{yes} in every cell \measured. Two rows deserve comment. At $k=2$, $N=10$
the measured ratio is $0.937$ against a ceiling of $1.125$: the coalition is
\emph{easier} to catch than one liar, because two statistics share the maximal
noncentrality and the maximum over two is larger than either. At the ends the
ratio is $1.000$ up to the bisection tolerance ($1.002$ at $k=1$, $1.003$ at
$k=9$), which is what the lemma predicts.

A second, sharper statement is measured in the same table: the chord test of
Remark~\ref{rem:chordtest} reaches power $1.000$ at $b_1$ for every interior $k$
while the node maximum reaches $0.793$--$0.954$ \measured. The information is
there; the node maximum is simply not the statistic that carries it. This is why
the dilution law of M1 is stated for the node maximum and not for the park: a
coalition of $k$ is not $k$ times less visible, it is $\sqrt{k(N-k)/(N-1)}$ times
less visible, and any test ordered by the chord sees that factor in full.

\begin{table*}[t]
\centering
\caption{The bound of Lemma~\ref{lem:bound} against the run (\texttt{theorem.measured}). $b_1$ is the measured minimum detectable lie for a lone liar at the declared risk $\alpha=10^{-3}$ and target power $0.9$; $\hat b_{\min}(k)$ is measured by bisection at the same power. ``pred'' is $(N-1)/\max(k,N-k)$, ``meas'' is $\hat b_{\min}(k)/b_1$, ``bound'' the inequality between them -- one column, because it is a ceiling and not an identity. ``chord'' is the power of the known-coalition test at the same $b_1$: a coalition is easier to catch than a node maximum suggests, never harder. ``flag on honest'' is the measured share of honest nodes flagged at $b_1$ -- the inversion of \S\ref{sec:inversion}, read from the same runs.}
\label{tab:measured}
\small
\begin{tabular}{@{}r r r r r r c r r r@{}}
\toprule
$N$ & $k$ & $b_1/\sigma$ & $\hat b_{\min}/\sigma$ & pred & meas & bound & power & chord & flag on honest \\
\midrule
10 & 1 & 5.420 & 5.432 & 1.000 & 1.002 & yes & 0.899 & 0.967 & 0.00036 \\
10 & 2 & 5.420 & 5.080 & 1.125 & 0.937 & yes & 0.954 & 1.000 & 0.00115 \\
10 & 3 & 5.420 & 5.238 & 1.286 & 0.966 & yes & 0.933 & 1.000 & 0.01134 \\
10 & 5 & 5.420 & 5.883 & 1.800 & 1.085 & yes & 0.793 & 1.000 & 0.50121 \\
10 & 6 & 5.420 & 5.607 & 1.500 & 1.035 & yes & 0.861 & 1.000 & 0.89433 \\
10 & 8 & 5.420 & 5.080 & 1.125 & 0.937 & yes & 0.952 & 1.000 & 0.99877 \\
10 & 9 & 5.420 & 5.438 & 1.000 & 1.003 & yes & 0.897 & 0.967 & 0.99975 \\
10 & 10 & 5.420 & $\infty$ & $\infty$ & $\infty$ & -- & 0.001 & -- & -- \\
30 & 1 & 5.513 & 5.511 & 1.000 & 1.000 & yes & 0.900 & 0.984 & 0.00022 \\
30 & 2 & 5.513 & 4.875 & 1.036 & 0.884 & yes & 0.983 & 1.000 & 0.00008 \\
30 & 5 & 5.513 & 4.465 & 1.160 & 0.810 & yes & 0.999 & 1.000 & 0.00010 \\
30 & 10 & 5.513 & 4.807 & 1.450 & 0.872 & yes & 0.992 & 1.000 & 0.01568 \\
30 & 15 & 5.513 & 5.354 & 1.933 & 0.971 & yes & 0.933 & 1.000 & 0.50019 \\
30 & 20 & 5.513 & 4.807 & 1.450 & 0.872 & yes & 0.992 & 1.000 & 0.98375 \\
30 & 25 & 5.513 & 4.460 & 1.160 & 0.809 & yes & 0.998 & 1.000 & 0.99982 \\
30 & 29 & 5.513 & 5.532 & 1.000 & 1.004 & yes & 0.898 & 0.983 & 0.99989 \\
30 & 30 & 5.513 & $\infty$ & $\infty$ & $\infty$ & -- & 0.001 & -- & -- \\
\bottomrule
\end{tabular}
\end{table*}


\section{Field validation}
\label{sec:field}

Everything above is a statement about a model. The run also puts the model on a
measured background: PVGIS \cite{pvgis} and NASA POWER \cite{nasa} irradiance and
temperature for a site
near Vladimir, four seasons, windows of 1, 24 and 168 hours, a physical
solar-geometry forward model turning weather into a frequency offset, and the
detector fed the resulting traces \measured. The background is not Gaussian: after
detrending, the excess kurtosis of the residual is between $-0.315$ and $1.689$
across the 32 measured cells, with 16 of them non-finite (degenerate phase cells)
\measured. The largest value, $1.689$, is in the July 24-hour offset cell -- a
heavy-tailed day, and exactly the cell the ROC is drawn on.

\begin{remark}[What survives non-Gaussianity]
The symmetry results do not care: Theorem~\ref{thm:kernel},
Lemma~\ref{lem:chord} and Lemma~\ref{lem:orientation} are statements about the
linear map $\PH$, and hold for any noise law. The quantitative ones do care, and
the run prices them: at $\alpha=10^{-3}$ the declared risk is paid as
$0.0008$--$0.00105$ of honest parks flagged \measured{}, i.e.\ the tail of the
measured background is close enough to Gaussian at the $3.9\sigma$ level that the
declared risk is not a fiction -- but it is a measured claim, not a theorem.
\end{remark}

Table~\ref{tab:field} is the collusion experiment on the real background: a
coalition of $k$ nodes lies by $b/\sigma=8$ in a park of $N$, 40 repetitions per
cell, seasons January and July. Two rows summarise the paper. At $N=10$, $k=9$:
detection $1.000$, localisation $0.000$ -- the detector fires every time and
names the honest node every time, which is
Corollary~\ref{cor:inversion} on weather rather than on white noise. At $N=10$,
$k=10$: detection $0.000$ -- a unanimous park is silent, which is
Theorem~\ref{thm:kernel} on weather. In between, localisation follows the
predicted side: \texttt{colluder} for $k<N/2$, \texttt{tie} at $k=N/2$,
\texttt{honest} above it \measured.

The window matters as much as the theory. At $T=24$~h the background is
non-stationary enough that the estimate of $\sigma$ absorbs the weather and the
measured colluder gain collapses to $\approx0.005$ instead of the predicted
$0.99$; at $T=168$~h the same code recovers the prediction \measured. We report
the 168-hour window for the theorem and both windows for the ROC, because a
theory validated only where it is convenient is not validated.


\begin{table}[t]
\centering
\caption{Field validation on the measured background (January, $T=168$~h, frequency offset, $b/\sigma=8$, 40 repetitions per cell). Detection is the share of repetitions crossing $\tau$; localisation is the share whose largest $|t|$ sits on a colluder. Past $k>N/2$ detection stays at $1.000$ while localisation collapses to $0.000$: the detector fires and names the honest, exactly as \S\ref{sec:inversion} requires. At $k=N$ there is nothing to detect and nothing is claimed.}
\label{tab:field}
\small
\begin{tabular}{@{}r r r r l@{}}
\toprule
$N$ & $k$ & detection & localisation & flagged side (pred) \\
\midrule
10 & 1 & 1.000 & 1.000 & colluder \\
10 & 2 & 1.000 & 1.000 & colluder \\
10 & 3 & 1.000 & 1.000 & colluder \\
10 & 5 & 1.000 & 0.575 & tie \\
10 & 6 & 1.000 & 0.025 & honest \\
10 & 8 & 1.000 & 0.000 & honest \\
10 & 9 & 1.000 & 0.000 & honest \\
10 & 10 & 0.000 & 1.000 & none \\
30 & 1 & 1.000 & 1.000 & colluder \\
30 & 2 & 1.000 & 1.000 & colluder \\
30 & 5 & 1.000 & 1.000 & colluder \\
30 & 10 & 1.000 & 1.000 & colluder \\
30 & 15 & 1.000 & 0.575 & tie \\
30 & 20 & 1.000 & 0.000 & honest \\
30 & 25 & 1.000 & 0.000 & honest \\
30 & 29 & 1.000 & 0.000 & honest \\
30 & 30 & 0.000 & 1.000 & none \\
\bottomrule
\end{tabular}
\end{table}


\section{Discussion: what is proved, and where the model is the weak link}
\label{sec:limits}

The results split cleanly by epistemic status, and the split is the most useful
thing in the paper.

\paragraph{Proved, model-free.} Lemma~\ref{lem:gauge}, Theorem~\ref{thm:kernel},
Lemma~\ref{lem:chord}, Proposition~\ref{prop:means},
Lemma~\ref{lem:orientation} and Lemma~\ref{lem:bound} are statements about the
linear map $\PH$ and about maxima of correlated normals. They do not assume
Gaussian noise, a particular estimator, a window length, or a threshold rule. If
the observation is the park and the park is read as contrasts, they hold.

\paragraph{Measured, model-conditional.} The power, the ROC, the coefficient
$c(N)$, the anchored prices and every number in Tables~\ref{tab:measured} and
\ref{tab:field} depend on \eqref{eq:h1} and on a background that is not Gaussian.
The excess kurtosis of the measured background reaches $1.689$ \measured{} in the
very cell on which the ROC is drawn; the declared risk is achieved there to within
$5\%$ of $\alpha$, which we report as an empirical fact about this background and
not as a guarantee about another.

\paragraph{Assumed, and therefore the weak link.} Four assumptions carry the
quantitative claims, and each has a known failure mode. (i) \emph{One bias per
coalition.} A coalition that splits into two opposite biases cancels in the
anchored mean and shrinks the chord; a clever coalition knows this, and
\S\ref{sec:conclusion} names it as the next statement to prove. (ii) \emph{A clock
that drifts, not one that jumps.} A timestamp step is a common-mode shift,
invisible to a frequency reading -- Theorem~\ref{thm:kernel} reached from the
other side. (iii) \emph{A stationary window.} At 24 hours the estimate of $\sigma$
absorbs weather and the measured colluder gain collapses; the theorem does not
fail, the estimator does. (iv) \emph{No substitution of history.} A node that
rewrites its own past is not covered by any statement here: the park compares
nodes, and a node that lies consistently in the past is, observationally, a node
with a different history -- the gauge again, one level up.

\paragraph{Relation to Kuramoto.} The classical reading of a park of oscillators
\cite{kuramoto1975} describes how coupling pulls frequencies into a common rhythm;
the order parameter measures \emph{how much} the park agrees. Our statement is
orthogonal and complementary: the order parameter is a function of the contrasts
and is therefore invariant to the common mode, which is exactly why a unanimous
lie cannot move it. Entrainment theory tells us when a park agrees; the gauge
symmetry tells us that agreement certifies nothing about the truth of the agreed
frequency. Together they give the operational reading of \S\ref{sec:anchor}:
agreement is free, and certification costs an anchor.

\paragraph{Relation to classical testing.} The GLRT is a vehicle here, not an
optimum: by Lemma~\ref{lem:chord} every admissible test sees the same chord, and
Remark~\ref{rem:chordtest} shows that a test ordered by the chord dominates one
ordered by the node maximum at interior $k$. Classical theory \cite{lehmann2005}
tells us the GLRT is efficient against a \emph{single} alternative; a coalition is
not one, and Table~\ref{tab:measured} is the price of that mismatch.

\section{Conclusion}
\label{sec:conclusion}

A park of oscillators can detect that it is split, and it cannot tell which side
is lying. That is not a limitation of the detector: it is the shape of the
observation, and it is provable. The provable part is small, sharp and free of
assumptions -- a projection, a chord, a kernel, a one-sided ceiling -- and it
constrains every future algorithm, including the ones that will replace ours. The
measured part is large, conditional and honest about it: a declared risk that is
paid, a ROC that does not improve with the size of the park, an inversion that
fires on weather exactly as it fires on white noise.

The practical reading is a trade, and it is the one we would defend in
deployment. If the threat is a few liars, use the plain GLRT and keep the park
flat: no node is privileged, and the floor against one liar is the best the park
can offer. If the threat is a majority, the park cannot help and no amount of data
will change that; the only move is to declare an anchor and pay $2.4\times$
($N=10$) to $4.3\times$ ($N=30$) worse sensitivity against the few \measured. An
anchor is not a better clock. It is an assumption, bought at a price, and the
price is now written down.

\paragraph{Next.} Three directions, in the order we intend to take them. First,
the mean-field park with memory (M3): the same quotient argument applies to a
coupled system with state, where the chord becomes a function of the coupling
matrix, the anchor becomes a boundary condition, and the fixed point of the
iteration is reached by a contraction \cite{banach1922}. Second, external validation
on a park whose truth is known independently of its own reports -- a laboratory
clock, or a second park on a different physical channel -- the only way to test
the orientation claim instead of assuming it. Third, the adversarial model: a
coalition that splits its bias to minimise the chord. The theorem says such a
coalition is doing the only thing that helps, and the size of the help is
computable; that is the next statement to prove.

\appendix

\section{Code}
\label{app:code}

The mathematics of this paper is \texttt{entrain\_infer/theorem.py} (436 lines):
\texttt{contrast}, \texttt{chord}, \texttt{chord\_ratio},
\texttt{colluder\_t\_mean}, \texttt{honest\_t\_mean}, \texttt{b\_eff\_factor},
\texttt{node\_gain}, \texttt{b\_min\_ratio}, \texttt{flagged\_side},
\texttt{two\_worlds}, \texttt{anchor\_reading}, \texttt{anchor\_statistic},
\texttt{anchor\_noncentrality}, \texttt{anchor\_b\_min}, \texttt{anchor\_power},
\texttt{anchor\_crossover}, \texttt{chord\_statistic}, \texttt{chord\_b\_min},
\texttt{chord\_power}. It is closed form throughout: no Monte Carlo, no seed, no
threshold. The detector and the estimator are \texttt{detector.py} and
\texttt{estimator.py}; the bench that produced every number here is
\texttt{bench.py} with the report in \texttt{report.py}.

\section{Tests}
\label{app:tests}

27 tests cover the theorem module (\texttt{tests/test\_theorem.py}, 335 lines) and
74 tests in total cover the package: the algebra of the projection, the chord
identity, the per-node means, the two worlds, the anchor, the bound and its
attainment, the kernel on the noise model and on the real background. The two that
matter most for \S\ref{sec:bound} are
\path{test_the_bound_is_attained_only_at_the_ends} (the ratio equals its ceiling
at $k=1$ and $k=N-1$, strictly below it in between) and
\path{test_a_coalition_at_the_bound_is_at_least_as_loud_as_one_liar} (power at the
bound: $0.895$ for one liar, $0.99$--$1.00$ for a coalition, $0.896$ at $k=N-1$).


\section{Reproducibility}
\label{app:repro}

Every number in this paper is read from \texttt{bench/results.json}, whose SHA-256
is
\begin{center}
\texttt{334b9939d4e055a204715d048fc4ae4b}\\
\texttt{c426ff40e3c65fd5b0bc5d4807a37683}
\end{center}
and which two independent full runs reproduce byte for byte. The two
pre-registrations, written before their runs and never edited afterwards, are
\texttt{bench/criteria.json} (\texttt{23dd1c85\ldots}) and
\texttt{bench/criteria\_m2.json} (\texttt{c8cfbad1\ldots}); the run prints both
hashes, and the report refuses to change them. The tables of this paper are
generated from that JSON by \texttt{paper/make\_tables.py} and contain no typed
numbers; \texttt{bench/REPORT.md} is regenerated deterministically from the same
file.

\small
\begin{thebibliography}{9}

\bibitem{kuramoto1975}
Y.~Kuramoto,
\newblock Self-entrainment of a population of coupled non-linear oscillators,
\newblock in \emph{International Symposium on Mathematical Problems of
Theoretical Physics}, Lecture Notes in Physics \textbf{39}, 420--422 (Springer,
1975).

\bibitem{lehmann2005}
E.~L. Lehmann and J.~P. Romano,
\newblock \emph{Testing Statistical Hypotheses}, 3rd ed.,
\newblock Springer, New York, 2005.

\bibitem{banach1922}
S.~Banach,
\newblock Sur les op\'erations dans les ensembles abstraits et leur application
aux \'equations int\'egrales,
\newblock \emph{Fundamenta Mathematicae} \textbf{3}, 133--181 (1922).

\bibitem{neyman1933}
J.~Neyman and E.~S. Pearson,
\newblock On the problem of the most efficient tests of statistical hypotheses,
\newblock \emph{Philosophical Transactions of the Royal Society A}
\textbf{231}, 289--337 (1933).

\bibitem{pvgis}
T.~Huld, R.~M\"uller and A.~Gambardella,
\newblock A new solar radiation database for estimating PV performance in Europe
and Africa,
\newblock \emph{Solar Energy} \textbf{86}, 1803--1815 (2012).
\newblock PVGIS-5 traces for the site used here, retrieved 2026.

\bibitem{nasa}
NASA Langley Research Center,
\newblock POWER Project Data Sets: meteorology and solar energy,
\newblock \emph{NASA Prediction of Worldwide Energy Resources}, 2021 data,
retrieved 2026.

\end{thebibliography}

\end{document}

